% GENERATED FILE. Edit canonical lessons, then run npm run generate:books.
% canonical-source-hash: fnv1a64:ab3f48847da6f985
% canonical-lessons: GE-C278-verhandeln, GE-C278-verhandlung, GE-C278-vertrag, GE-C278-unterzeichnen, GE-C278-konferenz

\chapter{Verhandlungen}
\label{ch:ge-verhandlungen}
\hlchaptermodality{278}

\begin{chapteropening}
\textbf{By the end of this chapter:} \emph{I can report talks and agreements: to negotiate, a negotiation, a contract, to sign and a conference.}

Complete the last lesson of chapter 278: I can report talks and agreements: to negotiate, a negotiation, a contract, to sign and a conference.
\end{chapteropening}

\section[verhandeln]{verhandeln — to negotiate}

\label{lesson:GE-C278-verhandeln}



\begin{quote}
Before the new one: say the German for to demonstrate, then the German for the crisis.
\end{quote}

\subsection*{You'll want to know: verhandeln}
\textbf{verhandeln} — \textquotedblleft{}to negotiate\textquotedblright{}.

\begin{quote}
Die beiden Firmen verhandelten eine ganze Woche.
\end{quote}

\begin{quote}
\emph{The two companies negotiated for a whole week.}
\end{quote}

\emph{Über etwas verhandeln} — to negotiate about something.

\begin{cousinweb}
\emph{verhandeln} is built on \emph{handeln} \textquotedblleft{}to deal\textquotedblright{}, from Old High German \emph{hantalōn}, which EtymDB derives from Proto-Germanic \emph{handuz} \textquotedblleft{}hand\textquotedblright{}.
\end{cousinweb}

\begin{grammarlens}[title={What you've built}]
One more word for telling, explaining and describing at B1.
\end{grammarlens}

\subsection*{Your turn}
\begin{itemize}
\raggedright
  \item \emph{Say it:} \emph{verhandeln}
  \item \emph{Say it:} \emph{verhandeln}, once more
  \item \emph{Say it:} a sentence of your own with \emph{verhandeln}
  \item \emph{Recall:} say the German for to demonstrate, then the German for the crisis, then say \emph{verhandeln} again
\end{itemize}

\begin{tcolorbox}[breakable,colback=teal!4,colframe=teal!35!black,title={Before you move on}]
What does verhandeln mean? (\textquotedblleft{}To negotiate\textquotedblright{}.) Say it once more.
\end{tcolorbox}
\section[die Verhandlung]{die Verhandlung — the negotiation; the hearing (in court)}

\label{lesson:GE-C278-verhandlung}



\begin{quote}
Before the new one: say the German for the crisis, then the German for to negotiate.
\end{quote}

\subsection*{You'll want to know: die Verhandlung}
\textbf{die Verhandlung} — \textquotedblleft{}the negotiation; the hearing (in court)\textquotedblright{}.

\begin{quote}
Die Verhandlung dauerte bis in die Nacht.
\end{quote}

\begin{quote}
\emph{The negotiation went on into the night.}
\end{quote}

Plural: \emph{die Verhandlungen}. In court, the hearing.

\begin{cousinweb}
\emph{Verhandlung} is built on \emph{verhandeln}, from \emph{handeln} \textquotedblleft{}to deal\textquotedblright{}, which EtymDB derives from Proto-Germanic \emph{handuz} \textquotedblleft{}hand\textquotedblright{}.
\end{cousinweb}

\begin{grammarlens}[title={What you've built}]
One more word for telling, explaining and describing at B1.
\end{grammarlens}

\subsection*{Your turn}
\begin{itemize}
\raggedright
  \item \emph{Say it:} \emph{die Verhandlung}
  \item \emph{Say it:} \emph{die Verhandlung}, once more
  \item \emph{Say it:} a sentence of your own with \emph{die Verhandlung}
  \item \emph{Recall:} say the German for the crisis, then the German for to negotiate, then say \emph{die Verhandlung} again
\end{itemize}

\begin{tcolorbox}[breakable,colback=teal!4,colframe=teal!35!black,title={Before you move on}]
What does die Verhandlung mean? (\textquotedblleft{}The negotiation; the hearing (in court)\textquotedblright{}.) Say it once more.
\end{tcolorbox}
\section[der Vertrag]{der Vertrag — the contract, the treaty}

\label{lesson:GE-C278-vertrag}



\begin{quote}
Before the new one: say the German for to negotiate, then the German for the negotiation.
\end{quote}

\subsection*{You'll want to know: der Vertrag}
\textbf{der Vertrag} — \textquotedblleft{}the contract, the treaty\textquotedblright{}.

\begin{quote}
Sie lasen den Vertrag genau.
\end{quote}

\begin{quote}
\emph{They read the contract carefully.}
\end{quote}

Plural: \emph{die Verträge}.

\begin{cousinweb}
\emph{Vertrag} is built on \emph{vertragen} \textquotedblleft{}to agree, to bear with\textquotedblright{}, from \emph{tragen} \textquotedblleft{}to carry\textquotedblright{}, Proto-Germanic \emph{draganą} \textquotedblleft{}to draw, to drag\textquotedblright{} (EtymDB).
\end{cousinweb}

\begin{grammarlens}[title={What you've built}]
One more word for telling, explaining and describing at B1.
\end{grammarlens}

\subsection*{Your turn}
\begin{itemize}
\raggedright
  \item \emph{Say it:} \emph{der Vertrag}
  \item \emph{Say it:} \emph{der Vertrag}, once more
  \item \emph{Say it:} a sentence of your own with \emph{der Vertrag}
  \item \emph{Recall:} say the German for to negotiate, then the German for the negotiation, then say \emph{der Vertrag} again
\end{itemize}

\begin{tcolorbox}[breakable,colback=teal!4,colframe=teal!35!black,title={Before you move on}]
What does der Vertrag mean? (\textquotedblleft{}The contract, the treaty\textquotedblright{}.) Say it once more.
\end{tcolorbox}
\section[unterzeichnen]{unterzeichnen — to sign (formally)}

\label{lesson:GE-C278-unterzeichnen}



\begin{quote}
Before the new one: say the German for the negotiation, then the German for the contract.
\end{quote}

\subsection*{You'll want to know: unterzeichnen}
\textbf{unterzeichnen} — \textquotedblleft{}to sign (formally)\textquotedblright{}.

\begin{quote}
Beide Seiten unterzeichneten den Vertrag.
\end{quote}

\begin{quote}
\emph{Both sides signed the treaty.}
\end{quote}

More formal than \emph{unterschreiben}. Not separable: \emph{hat unterzeichnet}.

\begin{cousinweb}
\emph{unterzeichnen} joins \emph{unter} \textquotedblleft{}under\textquotedblright{} and \emph{zeichnen} \textquotedblleft{}to mark\textquotedblright{}, from \emph{Zeichen} \textquotedblleft{}sign\textquotedblright{}, an inherited Germanic word; EtymDB lists no parent for the verb: to put one's mark under.
\end{cousinweb}

\begin{grammarlens}[title={What you've built}]
One more word for telling, explaining and describing at B1.
\end{grammarlens}

\subsection*{Your turn}
\begin{itemize}
\raggedright
  \item \emph{Say it:} \emph{unterzeichnen}
  \item \emph{Say it:} \emph{unterzeichnen}, once more
  \item \emph{Say it:} a sentence of your own with \emph{unterzeichnen}
  \item \emph{Recall:} say the German for the negotiation, then the German for the contract, then say \emph{unterzeichnen} again
\end{itemize}

\begin{tcolorbox}[breakable,colback=teal!4,colframe=teal!35!black,title={Before you move on}]
What does unterzeichnen mean? (\textquotedblleft{}To sign (formally)\textquotedblright{}.) Say it once more.
\end{tcolorbox}
\section[die Konferenz]{die Konferenz — the conference}

\label{lesson:GE-C278-konferenz}



\begin{quote}
Before the new one: say the German for the contract, then the German for to sign.
\end{quote}

\subsection*{You'll want to know: die Konferenz}
\textbf{die Konferenz} — \textquotedblleft{}the conference\textquotedblright{}.

\begin{quote}
Die Konferenz war im Rathaus.
\end{quote}

\begin{quote}
\emph{The conference was in the town hall.}
\end{quote}

Plural: \emph{die Konferenzen}. Stress at the end: \emph{Kon-fe-RENZ}.

\begin{cousinweb}
\emph{Konferenz} comes from Medieval Latin \emph{conferentia}, built on \emph{conferre} \textquotedblleft{}to bring together\textquotedblright{}, from \emph{con-} \textquotedblleft{}together\textquotedblright{} and \emph{ferre} \textquotedblleft{}to carry\textquotedblright{} (EtymDB).
\end{cousinweb}

\begin{grammarlens}[title={What you've built}]
One more word for telling, explaining and describing at B1.
\end{grammarlens}

\subsection*{Your turn}
\begin{itemize}
\raggedright
  \item \emph{Say it:} \emph{die Konferenz}
  \item \emph{Say it:} \emph{die Konferenz}, once more
  \item \emph{Say it:} a sentence of your own with \emph{die Konferenz}
  \item \emph{Recall:} say the German for the contract, then the German for to sign, then say \emph{die Konferenz} again
\end{itemize}

\begin{tcolorbox}[breakable,colback=teal!4,colframe=teal!35!black,title={Before you move on}]
What does die Konferenz mean? (\textquotedblleft{}The conference\textquotedblright{}.) Say it once more.
\end{tcolorbox}
